\documentclass[12pt,a4paper]{article}
\usepackage[utf8]{inputenc}
\usepackage[indonesian]{babel}
\usepackage[margin=1in]{geometry}
\usepackage{xcolor}
\usepackage{array}
\usepackage{tabularx}
\usepackage{tabularray}
\usepackage{graphicx}
\usepackage{fancyhdr}
\usepackage{hyperref}
\usepackage{multirow}
\usepackage{booktabs}
\usepackage{enumitem}
\usepackage{longtable}
\usepackage{textcomp}
\usepackage{amssymb}

% Color definitions
\definecolor{samsungblue}{RGB}{28, 106, 187}
\definecolor{lightgray}{RGB}{240, 240, 240}
\definecolor{darkgray}{RGB}{100, 100, 100}

% Header setup
\pagestyle{fancy}
\fancyhf{}
\rhead{\thepage}
\renewcommand{\headrulewidth}{0pt}

% Custom commands
\newcommand{\sectionbox}[1]{%
  \colorbox{samsungblue}{%
    \parbox{\textwidth}{%
      \textcolor{white}{\large\bfseries #1}%
    }%
  }%
}

\newcommand{\subsectionbox}[1]{%
  \textbf{\large #1}\\[0.3cm]
}

\newcommand{\greencheck}{$\checkmark$}
\newcommand{\redcross}{$\times$}
\newcommand{\warnsign}{$\triangle$}

\title{}
\author{}
\date{}

\begin{document}

% ============================================================================
% TITLE PAGE
% ============================================================================
\thispagestyle{empty}
\vspace*{1cm}

\begin{center}
  \huge\textbf{PROPOSAL IDE/CONCEPT PAPER}\\[0.5cm]
  \huge\textbf{SAMSUNG SOLVE FOR TOMORROW 2026}\\[2cm]
\end{center}

\normalsize

Tulis ide solusi yang ingin Anda ajukan dengan mengisi kolom-kolom di bawah secara lengkap dan spesifik! Jika ada pertanyaan, silakan hubungi di WhatsApp Enrico (6289677878563) atau di email: \texttt{sft2026@dibimbing.id} atau dapat bergabung ke community Telegram Samsung Solve for Tomorrow 2026.

\vspace{1.5cm}

\begin{tabularx}{\textwidth}{|l|X|}
  \hline
  \textbf{Nama Tim} & : \textit{Berkat ILLAHI} \\[0.3cm]
  \textbf{Nama Sekolah/Perguruan Tinggi} & : \textit{Universitas Darussalam Gontor (UNIDA Gontor)} \\[0.3cm]
  \textbf{Nama Ketua} & : \textit{FARREL GHOZY AFFIFUDIN} \\[0.3cm]
  \textbf{Nama Anggota 1} & : \textit{FATIH JAWWAD AL MUMTAZ} \\[0.3cm]
  \textbf{Nama Anggota 2} & : \textit{ZIA ZARIN ROSYIDI} \\[0.3cm]
  \textbf{Nama Anggota 3} & : \textit{MUHAMMAD RASYA NAUFAL KHADAFI} \\[0.3cm]
  \hline
\end{tabularx}

\newpage

% ============================================================================
% SECTION 1: JUDUL PROJEK/IDE
% ============================================================================
\sectionbox{1. JUDUL PROJEK/IDE}

\vspace{0.5cm}

\textbf{\Large Water Presence Monitoring System: Citizen Science + AI-Powered Remote Sensing untuk Deteksi dan Monitoring Real-Time Keberadaan Air Permukaan}

\vspace{1cm}

% ============================================================================
% SECTION 2: TEMA/BIDANG
% ============================================================================
\sectionbox{2. TEMA/BIDANG YANG DIPILIH}

\vspace{0.5cm}

\textbf{Sustainability \& Environment}

\vspace{1cm}

% ============================================================================
% SECTION 3: LATAR BELAKANG
% ============================================================================
\sectionbox{3. LATAR BELAKANG}

\vspace{0.5cm}

\subsectionbox{Konteks Masalah (Data \& Fakta)}

Indonesia sebagai negara dengan geografis kepulauan yang luas menghadapi tantangan serius dalam monitoring dan deteksi sumber daya air. Berdasarkan data dari Kementerian Pekerjaan Umum dan Perumahan Rakyat (PUPR), akses ke data water resources yang akurat dan real-time masih terbatas, terutama di daerah terpencil dan pasca-bencana.

Data menunjukkan:
\begin{itemize}
  \item Lebih dari 40\% daerah Indonesia masih kesulitan dalam pemetaan sumber daya air permukaan
  \item Survei manual untuk mendeteksi keberadaan air di area bencana memakan waktu 3-5 hari
  \item Satelit konvensional (Landsat) memiliki resolusi 30m, tidak cukup untuk area lokal
  \item Biaya survey ground-based berkisar 5-10 juta rupiah per lokasi
  \item Tidak ada platform yang mengintegrasikan citizen observations dengan data satelit secara real-time
  \item Citra satelit sering terhalang cloud cover (30-50\% tidak dapat digunakan), menyebabkan keterlambatan data
\end{itemize}

\subsectionbox{Riset Pengguna (Persona \& Kutipan Wawancara)}

\textbf{Persona 1: Disaster Response Manager}
\begin{itemize}
  \item Profil: Koordinator bencana di Badan Nasional Penanggulangan Bencana (BNPB)
  \item Tantangan: ``Kami butuh informasi cepat tentang area yang tergenang untuk prioritas evakuasi, tapi data satelit lambat dan mahal''
  \item Kebutuhan: Pemetaan real-time area tergenang dalam hitungan menit, bukan hari
\end{itemize}

\textbf{Persona 2: Water Resource Manager}
\begin{itemize}
  \item Profil: Petugas dinas pertanian di kabupaten rural
  \item Tantangan: ``Susah monitor ketersediaan air di area pertanian kami yang luas tanpa data akurat''
  \item Kebutuhan: Indikator ketersediaan air permukaan per musim untuk perencanaan irigasi
\end{itemize}

\textbf{Persona 3: Environmental Scientist}
\begin{itemize}
  \item Profil: Peneliti di universitas yang fokus pada environmental monitoring
  \item Tantangan: ``Data satellite mahal dan jarang tersedia cloud-free, perlu alternative untuk baseline assessment''
  \item Kebutuhan: Tool yang accessible dan dapat di-crowdsource untuk citizen science
\end{itemize}

\subsectionbox{Pernyataan Masalah Inti (How Might We...)}

\textbf{``How might we create an accessible, real-time water presence detection system that combines smartphone observations with satellite validation and AI analysis, enabling rapid assessment for disaster response, agricultural planning, and environmental monitoring?''}

\vspace{1cm}

% ============================================================================
% SECTION 4: TARGET PENGGUNA
% ============================================================================
\sectionbox{4. TARGET PENGGUNA}

\vspace{0.5cm}

\begin{tabularx}{\textwidth}{|l|X|}
  \hline
  \textbf{Primary Users} & Petugas disaster response (BNPB), Water resource managers (dinas pertanian/PUPR), Environmental scientists \\[0.3cm]
  \textbf{Secondary Users} & General public, NGO untuk environmental monitoring, Mahasiswa untuk citizen science \\[0.3cm]
  \textbf{Geographic Scope} & Indonesia, khususnya daerah rural \& disaster-prone areas \\[0.3cm]
  \textbf{Usia/Demografis} & 20-55 tahun, tech-literate, dengan access ke smartphone \\[0.3cm]
  \textbf{Volume Pengguna} & MVP: 100-500 pengguna; Scale-up: 5,000+ pengguna dalam 1 tahun \\
  \hline
\end{tabularx}

\vspace{1cm}

% ============================================================================
% SECTION 5: DESKRIPSI IDE
% ============================================================================
\sectionbox{5. DESKRIPSI IDE}

\vspace{0.5cm}

\subsectionbox{A. Konsep Solusi: Penjelasan Umum Ide}

\textbf{Water Presence Monitoring System} adalah platform terintegrasi yang menjadikan analisis satelit multi-sumber sebagai inti deteksi, dengan AI sebagai analis yang merekap dan menyimpulkan data secara cerdas:

\begin{enumerate}
  \item \textbf{Multi-Source Satellite Analysis (Google Earth Engine)}: Deteksi air permukaan menggunakan kombinasi data satelit:
  \begin{itemize}
    \item \textbf{Sentinel-1 SAR (Radar)}: Water detection via backscatter analysis — \textbf{tembus awan}, cocok untuk Indonesia
    \item \textbf{Sentinel-2 (Optis)}: NDWI (Normalized Difference Water Index) sebagai data pendukung saat awan rendah
    \item \textbf{CHIRPS}: Data curah hujan 7 hari terakhir untuk konteks kebasahan
    \item \textbf{OpenLandMap}: Klasifikasi jenis tanah untuk konteks tambahan
    \item \textbf{SRTM/DEM}: Data elevasi untuk analisis topografi
  \end{itemize}
  \item \textbf{AI Analysis (Gemini 2.0 Flash)}: Menganalisis \textbf{seluruh data satelit} yang telah dikumpulkan GEE. Gemini bertugas merangkum, menginterpretasi, dan memberikan kesimpulan akurat tentang keberadaan air beserta confidence score.
  \item \textbf{Citizen Observation}: Cukup dengan GPS — user memilih lokasi di peta atau menggunakan GPS otomatis. Proses deteksi air sepenuhnya dari data satelit.
\end{enumerate}

\textbf{Keluaran sistem}: Confidence score 0-100\% yang dihasilkan dari analisis Gemini terhadap data satelit, disertai dengan penjelasan natural language, breakdown data satelit per sumber, visualisasi di peta interaktif Indonesia (choropleth per provinsi), dan rekomendasi.

\textbf{Waktu processing}: 15-30 detik per observation (GEE query time dominan)

\subsectionbox{B. Rancangan Awal: User Flow dengan Penjelasan}

\vspace{0.3cm}

\textbf{\underline{User Flow Diagram:}}

\begin{verbatim}
User klik "Submit Observation"
    ↓
Pilih lokasi via GPS (auto-detect / klik peta)
    ↓
Backend API → POST /api/observations
    ↓
GOOGLE EARTH ENGINE — MULTI-SOURCE (15-30 detik)
    → Sentinel-1 SAR: water mask (tembus awan)
    → Sentinel-2 NDWI: jika cloud < 20%
    → CHIRPS: curah hujan 7 hari
    → OpenLandMap: tipe tanah
    → SRTM: elevasi & terrain
    ↓
Semua data → GEMINI AI ANALYST
    → Menganalisis SAR + NDWI + rainfall + soil + elevasi
    → Confidence score (0-100%)
    → Natural language assessment + rekomendasi
    ↓
TAMPILKAN KE USER
    → Confidence gauge + verdict
    → AI reasoning (natural language)
    → Peta Indonesia choropleth
    → Satellite data breakdown per source
    → Rekomendasi tindakan
\end{verbatim}

\vspace{0.3cm}

\textbf{\underline{Hubungan dengan Riset Pengguna:}}

\begin{itemize}
  \item \textbf{Disaster Response Manager}: Mendapat alert + peta real-time dalam 15 detik untuk rapid assessment
  \item \textbf{Water Resource Manager}: Dapat track ketersediaan air per lokasi dengan confidence metrics untuk planning
  \item \textbf{Environmental Scientist}: Dapat participate dalam citizen science dengan transparent methodology
\end{itemize}

\subsectionbox{C. Fitur Utama Aplikasi}

\begin{enumerate}
  \item \textbf{Observation Submission}
  \begin{itemize}
    \item GPS auto-detection dengan opsi manual adjustment via klik pada map
    \item Opsional: time window selector untuk satellite data lookback
  \end{itemize}

  \item \textbf{Results Dashboard (3-Panel Layout)}
  \begin{itemize}
    \item Panel 1: AI Analysis summary — confidence gauge 0-100\%, natural language assessment dari Gemini, water presence verdict, risk indicators
    \item Panel 3: Interactive map — Peta Indonesia choropleth (warna per provinsi berdasarkan water index), marker lokasi, satellite overlay (SAR water mask)
    \item Panel 4: Satellite data breakdown — Detail dari setiap sumber satelit (SAR, NDWI, CHIRPS, soil, elevation) yang digunakan Gemini untuk analisis
  \end{itemize}

  \item \textbf{Explore/Map View}
  \begin{itemize}
    \item Interactive map (Leaflet.js) dengan \textbf{choropleth layer Indonesia}: setiap provinsi diwarnai berdasarkan water index dari satelit
    \item Heatmap overlay dari observasi crowdsourced
    \item Layer kontrol: pilih provinsi, ubah basemap, toggle satellite overlay
    \item Filters: date range, confidence threshold, water presence level
    \item Regional statistics \& trends (time-series chart per province/district)
    \item Klik provinsi untuk melihat detail water index dan observasi di dalamnya
  \end{itemize}

  \item \textbf{Methodology \& About}
  \begin{itemize}
    \item How the system works (transparent explanation with diagrams)
    \item Limitations \& disclaimers (honest about 60-80\% accuracy)
    \item Data sources attribution (Gemini, Sentinel-2, BMKG, OpenStreetMap)
    \item FAQ untuk pengguna umum
  \end{itemize}
\end{enumerate}

\subsectionbox{D. Arsitektur Sistem}

Sistem dirancang dengan arsitektur microservices yang terpisah antara frontend, backend, dan worker processing:

\begin{verbatim}
DEPLOYMENT ARCHITECTURE

Frontend (Vercel) ←→ Backend (Railway)
    React 19 + Vite       Bun + ElysiaJS
    Leaflet.js Maps        REST API
    Zustand + React Query  Gemini Integration
                           GEE Proxy
                              ↓
Backend panggil 3 service:
  → MongoDB Atlas (observations, results)
  → Python Worker (GEE multi-source)
      → Sentinel-1 SAR
      → Sentinel-2 NDWI
      → CHIRPS rainfall
      → OpenLandMap soil
      → SRTM elevation
  → Google Earth Engine API
\end{verbatim}

\textbf{Komponen Arsitektur:}

\begin{itemize}
  \item \textbf{Frontend}: React 19 + Vite dengan PWA capability. State management: Zustand + React Query. Map integration: Leaflet.js dengan plugin choropleth, GeoJSON overlay batas administrasi Indonesia, dan heatmap.
  \item \textbf{Backend API}: Bun runtime + ElysiaJS framework. REST API endpoints sederhana untuk observations dan results. Middleware: CORS, rate limiter, JWT validation, error handling.
  \item \textbf{Python Worker}: Microservice untuk Google Earth Engine queries multi-sumber. Mengumpulkan data dari Sentinel-1 SAR, Sentinel-2, CHIRPS, OpenLandMap, dan SRTM secara paralel.
  \item \textbf{Gemini Integration}: Backend mengirimkan seluruh data satelit (terstruktur) ke Gemini 2.0 Flash untuk analisis dan summarization. Gemini menghasilkan confidence score + natural language assessment.
  \item \textbf{Caching Strategy}: Data satelit (SAR, NDWI, CHIRPS) di-cache 7 hari. Soil dan elevation di-cache 30 hari. Analysis results tidak pernah expire.
\end{itemize}

\subsectionbox{E. Rancangan Database}

\begin{tabularx}{\textwidth}{|l|X|}
  \hline
  \textbf{MongoDB (Single Database)} & Satu database untuk semua data: \\[0.2cm]
  & - \texttt{observations}: userId, lokasi (lat/lng), timestamp, status (processing/completed/error) \\[0.2cm]
  & - \texttt{satellite\_data}: observationId, sar\_water\_mask, ndwi\_value, chirps\_rainfall, soil\_type, elevation \\[0.2cm]
  & - \texttt{gemini\_analysis}: observationId, confidence\_score, water\_presence\_assessment, natural\_language\_summary, recommendations \\[0.2cm]
  & - \texttt{regional\_index}: province, water\_index, last\_updated, observation\_count \\
  \hline
\end{tabularx}

\textbf{Data untuk peta}: Regional water index dihitung periodik dari agregasi data satelit per provinsi, disimpan sebagai dokumen terpisah untuk akses cepat visualisasi choropleth.

\subsectionbox{F. API Endpoints}

\begin{tabularx}{\textwidth}{|l|l|X|}
  \hline
  \textbf{Method} & \textbf{Endpoint} & \textbf{Deskripsi} \\
  \hline
  POST & /api/v1/observations & Submit observasi (GPS location) \\
  GET & /api/v1/observations/:id & Detail observasi + hasil analisis \\
  GET & /api/v1/observations/:id/satellite-data & Data satelit mentah per sumber \\
  GET & /api/v1/map/indonesia & GeoJSON peta Indonesia choropleth (water index per provinsi) \\
  GET & /api/v1/regions/:province/stats & Statistik regional \\
  GET & /api/v1/regions/:province/trends & Time-series data regional \\
  GET & /api/v1/health & System health check \\
  \hline
\end{tabularx}

\subsectionbox{G. Rancangan Teknis Singkat}

\begin{tabularx}{\textwidth}{|l|X|}
  \hline
  \textbf{Frontend} & React 19 + Vite + Tailwind CSS + Leaflet.js + Zustand + React Query \\[0.2cm]
  \textbf{Backend} & Bun runtime + ElysiaJS framework \\[0.2cm]
  \textbf{AI/Analyst} & Google Gemini 2.0 Flash (analisis data satelit \& summarization) \\[0.2cm]
  \textbf{Satellite Radar} & Sentinel-1 SAR via GEE — water detection (tembus awan) \\[0.2cm]
  \textbf{Satellite Optis} & Sentinel-2 via GEE — NDWI (jika awan rendah) \\[0.2cm]
  \textbf{Rainfall} & CHIRPS via GEE — data curah hujan \\[0.2cm]
  \textbf{Soil \& Elevation} & OpenLandMap + SRTM via GEE \\[0.2cm]
  \textbf{Geospatial} & Google Earth Engine (platform integrasi semua data satelit) \\[0.2cm]
  \textbf{Database} & MongoDB Atlas (observations \& analysis results) \\[0.2cm]
  \textbf{Caching} & Redis / in-memory cache untuk data satelit (7 hari TTL) \\[0.2cm]
  \textbf{Monitoring} & Sentry (error tracking) + Winston (logging) \\[0.2cm]
  \textbf{Deployment} & Vercel (frontend) + Railway (backend + Python worker) \\
  \hline
\end{tabularx}

\vspace{1cm}

% ============================================================================
% SECTION 6: TUJUAN
% ============================================================================
\sectionbox{6. TUJUAN}

\vspace{0.5cm}

\begin{enumerate}
  \item \textbf{Percepatan Deteksi}: Mengurangi waktu deteksi keberadaan air dari 3-5 hari (survey manual) menjadi 10-15 detik (sistem otomatis)

  \item \textbf{Aksesibilitas}: Membangun tool yang accessible via smartphone saja, tanpa butuh expertise khusus atau peralatan mahal

  \item \textbf{Akurasi}: Mencapai confidence score yang reliable (F1-score 0.65-0.80) melalui validasi multi-sumber (SAR radar + satellite optis + weather)

  \item \textbf{Skalabilitas}: Menghasilkan platform yang dapat scale ke ribuan pengguna dan observasi per hari dengan serverless architecture

  \item \textbf{Transparansi}: Provide confidence metrics \& methodology disclosure sehingga pengguna understand keterbatasan sistem

  \item \textbf{Dampak Sosial}: Mendukung disaster response, agricultural planning, dan environmental monitoring di Indonesia
\end{enumerate}

\vspace{1cm}

% ============================================================================
% SECTION 7: INOVASI
% ============================================================================
\sectionbox{7. INOVASI}

\vspace{0.5cm}

\subsectionbox{Keunikan Ide dibanding Solusi Existing:}

\begin{enumerate}
  \item \textbf{Multi-Source Satellite Analysis (Bukan Satu Sumber)}
  \begin{itemize}
    \item Existing tools: Analisis berdasarkan foto saja, atau satu sumber satelit
    \item Kami: Kombinasi 5+ sumber data satelit (SAR, optis, rainfall, soil, elevasi) via GEE — lebih akurat dan objektif
  \end{itemize}

  \item \textbf{SAR Radar untuk Indonesia (Tembus Awan)}
  \begin{itemize}
    \item Existing: Pakai optis (Landsat/Sentinel-2) yang 70-90\% tertutup awan di Indonesia
    \item Kami: Sentinel-1 SAR Radar — deteksi air permukaan \textbf{walau mendung}, cocok untuk iklim tropis
  \end{itemize}

  \item \textbf{AI sebagai Analyst, Bukan Comparator}
  \begin{itemize}
    \item Existing: AI untuk analisis foto atau sebagai pembanding sumber data
    \item Kami: Gemini menganalisis \textbf{seluruh dataset satelit} dan memberikan kesimpulan natural language yang komprehensif
  \end{itemize}

  \item \textbf{Peta Tematik Indonesia Otomatis}
  \begin{itemize}
    \item Existing: Peta air statis atau butuh update manual
    \item Kami: Peta choropleth Indonesia yang terupdate otomatis dari data satelit — warna per provinsi berdasarkan water index
  \end{itemize}

  \item \textbf{Cost-Effective Solution}
  \begin{itemize}
    \item Existing: Professional survey: 5-10M per lokasi, Satellite imagery: 1-5M per order
    \item Kami: GEE free tier untuk akses semua data satelit, Gemini untuk analisis → accessible untuk community-level
  \end{itemize}
\end{enumerate}

\vspace{1cm}

% ============================================================================
% SECTION 8: DAMPAK IDE/PROYEK
% ============================================================================
\sectionbox{8. DAMPAK IDE/PROYEK}

\vspace{0.5cm}

\subsectionbox{Dampak Langsung (Direct Impact):}

\begin{enumerate}
  \item \textbf{Disaster Response} (BNPB, disaster managers)
  \begin{itemize}
    \item Reduce assessment time dari hari menjadi menit
    \item Better spatial mapping untuk prioritasi evakuasi
    \item Estimated: Bisa accelerate response 2-3 hari lebih cepat per event bencana
  \end{itemize}

  \item \textbf{Agricultural Planning} (Dinas Pertanian, farmers)
  \begin{itemize}
    \item Real-time indicator ketersediaan air untuk irrigation planning
    \item Reduce crop loss dari drought via early detection
    \item Estimated: +10-15\% yield improvement di dry seasons
  \end{itemize}

  \item \textbf{Environmental Monitoring} (Researchers, NGOs)
  \begin{itemize}
    \item Baseline data collection untuk water resources assessment
    \item Crowdsourced validation untuk environmental projects
    \item Estimated: 5x lebih banyak data points dibanding traditional surveys
  \end{itemize}
\end{enumerate}

\subsectionbox{Dampak SDGs (Sustainable Development Goals):}

\begin{itemize}
  \item \textbf{SDG 6 (Clean Water \& Sanitation)}: Improve monitoring \& management of water resources
  \item \textbf{SDG 13 (Climate Action)}: Support disaster preparedness \& response dengan data real-time
  \item \textbf{SDG 15 (Life on Land)}: Environmental monitoring untuk conservation dan ecosystem management
\end{itemize}

\subsectionbox{Dampak Jangka Panjang:}

\begin{itemize}
  \item Foundation untuk integrated water resource management system di Indonesia
  \item Scalable model untuk citizen science platform di bidang environment yang bisa diadopsi negara lain
  \item Potential integration dengan government agencies (BNPB, PUPR, Kementan)
  \item Dataset crowdsourced yang dapat digunakan untuk penelitian dan kebijakan publik
\end{itemize}

\vspace{1cm}

% ============================================================================
% SECTION 9: ASPEK STEM
% ============================================================================
\sectionbox{9. ASPEK STEM}

\vspace{0.5cm}

\begin{tabularx}{\textwidth}{|l|X|}
  \hline
  \textbf{Science} &
    Aplikasi ilmu Hydrology \& Remote Sensing: NDWI (Normalized Difference Water Index) = (Green - NIR) / (Green + NIR) menggunakan Sentinel-2 band B3 (Green) dan B8 (NIR) untuk deteksi water bodies dari multispectral imagery. Threshold NDWI $>$ 0.3 mengindikasikan keberadaan air. Environmental science principles untuk data interpretation dan korelasi dengan data cuaca. \\[0.3cm]
  \hline
  \textbf{Technology} &
    Full-stack modern development: React 19 + Vite (frontend), Bun/ElysiaJS (backend), Python microservice (satellite processing). Cloud infrastructure: Vercel, Railway, MongoDB Atlas, Supabase.     APIs: Google Gemini 2.0 Flash, Google Earth Engine, CHIRPS. \\[0.3cm]
  \hline
  \textbf{Engineering} &
    Software engineering: REST API design, database architecture (MongoDB), microservices pattern (Python worker terpisah untuk GEE processing). System integration: menghubungkan 3+ external APIs secara seamless. UI/UX design untuk accessibility di mobile dengan validasi form real-time. \\[0.3cm]
  \hline
  \textbf{Mathematics} &
    Confidence score algorithm: weighted averaging dari multiple data sources. Formula: confidence = (0.4 \texttimes{} gemini\_conf) + (0.4 \texttimes{} satellite\_conf) + (0.2 \texttimes{} agreement\_score). NDWI formula dengan normalized difference. Statistical analysis untuk accuracy metrics (precision, recall, F1-score). Time-series analysis untuk regional trends dan anomaly detection. \\
  \hline
\end{tabularx}

\vspace{1cm}

% ============================================================================
% SECTION 10: KONTRIBUSI AI
% ============================================================================
\sectionbox{10. KONTRIBUSI AI}

\vspace{0.5cm}

\subsectionbox{Peran AI dalam Sistem:}

\begin{enumerate}
  \item \textbf{Gemini 2.0 Flash sebagai AI Analyst (Satellite Data Analysis)}
  \begin{itemize}
    \item Input: Seluruh data satelit yang telah dikumpulkan GEE dalam format terstruktur:
    \begin{itemize}
      \item Sentinel-1 SAR: water mask, backscatter values, persentase area basah
      \item Sentinel-2: NDWI value, cloud cover
      \item CHIRPS: curah hujan 7 hari terakhir, tren
      \item OpenLandMap: soil type classification
      \item SRTM: elevation, slope
      \item Metadata: timestamp, lokasi (provinsi/kabupaten)
    \end{itemize}
    \item Processing: Menganalisis semua data, mengidentifikasi pola, menilai konsistensi antar sumber, memberikan kesimpulan:
    \begin{itemize}
      \item Water presence assessment (definitive/probable/possible/unlikely)
      \item Confidence score 0-100\% berdasarkan kekuatan bukti dari data satelit
      \item Natural language explanation untuk setiap faktor yang mempengaruhi
      \item Rekomendasi tindakan (verifikasi lapangan / aman / pantau lanjut)
      \item Anomaly detection (ketidaksesuaian antar sumber data)
    \end{itemize}
    \item Output: JSON structured + natural language string
    \item Cost: ~USD 0.001-0.003 per analysis (tergantung length input)
    \item Latency: 3-8 detik per call
  \end{itemize}
\end{enumerate}

\subsectionbox{Flow AI Analysis:}

\begin{verbatim}
STEP 1: GEE mengumpulkan data satelit
  → SAR: {waterPercentage: 34, confidence: "high"}
  → NDWI: {value: 0.42, cloudCover: 15%}
  → CHIRPS: {rainfall7day: 120mm, trend: "increasing"}
  → Soil: {type: "clay loam"}
  → Elevation: {value: 45m, terrain: "flat"}

STEP 2: Kirim ke Gemini sebagai structured prompt
  → Format: JSON data satelit + metadata
  → System prompt: "You are a hydrology analyst..."
  → User prompt: "Analyze this satellite data..."

STEP 3: Gemini merespon
  → confidence: 78
  → verdict: "probable water presence"
  → reasoning: "SAR shows 34% water coverage..."
  → recommendations: ["Verify on site if critical"]
  → anomalies: ["NDWI limited by 15% cloud cover"]
\end{verbatim}

\subsectionbox{Mengapa AI Critical:}

\begin{itemize}
  \item Data satelit multi-sumber sulit diinterpretasi manual (5+ dataset berbeda)
  \item Gemini mampu meng-cross-reference data SAR, optis, rainfall, dan geospasial dalam satu konteks
  \item Natural language output lebih informatif daripada confidence score angka saja
  \item Tanpa AI: butuh ahli remote sensing untuk interpretasi data ini
  \item Dengan AI: masyarakat umum bisa mendapat analisis setara ahli dalam 5 detik
\end{itemize}

\subsectionbox{Detail Prompt Engineering untuk Gemini:}

\begin{verbatim}
System Prompt:
"You are a hydrology and remote sensing analyst specializing
in water presence detection in tropical environments. You
will receive multi-source satellite data for a specific
location. Analyze ALL data sources and respond in valid JSON."

Input Data (structured):
{
  "location": {"lat": -7.25, "lng": 112.75, "province": "Jatim"},
  "sar_analysis": {
    "water_percentage": 34.2,
    "backscatter_mean": -18.5,
    "confidence": "high"
  },
  "ndwi": {"value": 0.42, "cloud_cover": 15, "available": true},
  "rainfall": {"total_7day_mm": 120, "trend": "increasing"},
  "soil": {"type": "clay loam", "drainage": "moderate"},
  "elevation": {"meters": 45, "terrain": "flat"},
  "timestamp": "2026-05-23T10:30:00Z"
}

Response JSON:
{
  "confidence": 78,
  "waterPresence": "probable",
  "reasoning": "SAR shows 34% water coverage ...",
  "contributingFactors": ["rainfall 120mm"],
  "anomalies": [],
  "recommendations": ["Site verification recommended"]
}
\end{verbatim}

\vspace{1cm}

% ============================================================================
% SECTION 11: SUSTAINABILITY
% ============================================================================
\sectionbox{11. SUSTAINABILITY}

\vspace{0.5cm}

\subsectionbox{Strategi Keberlanjutan (Sustainability Strategy):}

\begin{enumerate}
  \item \textbf{Business Model}
  \begin{itemize}
    \item \textbf{Freemium}: Observations \& basic results gratis untuk public users
    \item \textbf{Government Partnerships}: Premium subscription untuk BNPB, Dinas Pertanian, PUPR (dengan dedicated support, advanced analytics, API access)
    \item \textbf{NGO Programs}: Discounted rates untuk research \& conservation projects
    \item \textbf{Estimated revenue}: IDR 500M - 1B per tahun di year 2-3
  \end{itemize}

  \item \textbf{Technical Sustainability}
  \begin{itemize}
    \item Serverless architecture (Vercel, Railway) auto-scaling, minimal ops overhead
    \item Open-source components (React, Leaflet, Tailwind) dengan strong community support
    \item API-first design memudahkan integrasi dengan third-party tools dan government systems
    \item Regular updates based on user feedback \& technology trends
    \item Python worker dapat di-extend untuk additional data sources di masa depan
  \end{itemize}

  \item \textbf{Community Engagement}
  \begin{itemize}
    \item Citizen science program: reward active contributors (badges, leaderboards, contribution recognition)
    \item Educational outreach: training workshops untuk disaster managers, farmers, dan mahasiswa
    \item Partnership dengan universitas untuk research validation dan dataset development
    \item Transparent data sharing (CC-BY license) untuk credibility dan academic use
  \end{itemize}

  \item \textbf{Impact Measurement}
  \begin{itemize}
    \item Monthly KPIs: observation count, unique users, accuracy metrics, response time
    \item Quarterly user satisfaction surveys untuk continuous improvement
    \item Annual impact report: lives helped, disaster response time reduced, area monitored, publications
  \end{itemize}

  \item \textbf{Data Governance}
  \begin{itemize}
    \item Open data: All observations available via public API dengan CC-BY license
    \item Privacy-first: No personal data collected beyond email (optional) untuk authentication
    \item Data retention: 5 years on MongoDB, archival to cold storage (S3 Glacier) setelahnya
    \item Data satelit di-cache 7 hari untuk menghemat kuota GEE
  \end{itemize}
\end{enumerate}

\vspace{1cm}

% ============================================================================
% PAGE 2: INFORMASI TAMBAHAN
% ============================================================================
\newpage

\sectionbox{INFORMASI TAMBAHAN}

\vspace{0.5cm}

\subsectionbox{A. Timeline Pengembangan (Development Roadmap)}

\textbf{Week 1: Foundation \& Setup (25-30 jam)}
\begin{itemize}
  \item Setup environment: GitHub repo, monorepo structure (frontend/ + backend/), Docker Compose (MongoDB, PostgreSQL, Redis)
  \item Frontend: React + Vite project, Tailwind CSS, Leaflet, routing, layout components, halaman Home
  \item Backend: ElysiaJS server skeleton, database models (Mongoose), REST routes (CRUD observations), error handling middleware
  \item Testing: Frontend $\rightarrow$ Backend integration, mobile responsive, Postman/curl verification
\end{itemize}

\textbf{Week 2: AI \& Satellite Integration (30-35 jam)}
\begin{itemize}
  \item Gemini API: geminiService.ts wrapper untuk analisis data satelit, error handling + retries
  \item Google Earth Engine: Service account setup, Python microservice, NDWI calculation function, Redis caching (7-day TTL)
  \item Data Flow Integration: 3 parallel jobs (Gemini + GEE + BMKG), job completion check, comparison trigger
  \item Tantangan: Gemini response format (force JSON), GEE learning curve, cloud cover filtering
\end{itemize}

\textbf{Week 3: Analysis Engine \& Results Page (30-35 jam)}
\begin{itemize}
  \item Comparison Engine: calculateConfidence() function, weighted scoring algorithm, verdict generator, anomaly detection
  \item Results Page UI: 3-panel layout (analysis summary, map, satellite breakdown), polling dengan React Query
  \item Explore/Map View: Interactive map (Leaflet), heatmap overlay, filters (date range, confidence threshold)
  \item Backend Analytics: /map/heatmap endpoint (GeoJSON), /regions/:region/stats, regional aggregation
\end{itemize}

\textbf{Week 4: Testing, Polish \& Launch Prep (30-35 jam)}
\begin{itemize}
  \item Validation Testing: 5-10 real observations di berbagai environment (urban, agricultural, water body, dry area, wet area)
  \item Bug Fixes: Critical bugs, performance optimization, caching, database indices
  \item Documentation: README, API docs, methodology, limitations, setup guide, screenshots
  \item Deployment: Vercel (frontend), Railway (backend), production testing, rollback plan
\end{itemize}

\begin{tabularx}{\textwidth}{|l|l|X|}
  \hline
  \textbf{Week} & \textbf{Total Hours} & \textbf{Key Deliverable} \\
  \hline
  1 & 25-30 & Functional form submit $\rightarrow$ DB save $\rightarrow$ return ID \\
  2 & 30-35 & End-to-end: submit $\rightarrow$ Gemini + GEE $\rightarrow$ stored \\
  3 & 30-35 & Complete results page + explore map + analysis engine \\
  4 & 30-35 & Production-ready MVP, validated, deployed \& tested \\
  \textbf{Total} & \textbf{115-135} & \textbf{$\sim$4 weeks full-time equivalent} \\
  \hline
\end{tabularx}

\vspace{0.5cm}

\subsectionbox{B. Tech Stack Summary}

\begin{tabularx}{\textwidth}{|l|X|}
  \hline
  \textbf{Frontend} & React 19 + Vite + Tailwind CSS + Leaflet.js + Zustand + React Query \\[0.2cm]
  \textbf{Backend} & Bun + ElysiaJS dengan middleware CORS, rate limiter, JWT \\[0.2cm]
  \textbf{Python Worker} & FastAPI microservice untuk GEE dan BMKG processing \\[0.2cm]
  \textbf{Database} & MongoDB Atlas (observations \& analysis) + PostgreSQL Supabase (time-series + PostGIS) + Redis Upstash (cache \& queue) \\[0.2cm]
  \textbf{AI/APIs} & Gemini 2.0 Flash, Google Earth Engine, BMKG API (fallback: OpenWeather) \\[0.2cm]
  \textbf{Queue} & Bull (Redis-backed, 5 job types, retry 2x, exponential backoff, dead-letter queue) \\[0.2cm]
  \textbf{Cache Layer} & Redis / in-memory cache untuk data satelit \\[0.2cm]
  \textbf{Deployment} & Vercel (frontend) + Railway (backend + Python worker) \\[0.2cm]
  \textbf{Monitoring} & Sentry (error tracking) + Winston (structured logging) + Prometheus (/metrics endpoint) \\
  \hline
\end{tabularx}

\vspace{0.5cm}

\subsectionbox{C. Data Processing Pipeline \& Confidence Scoring}

\textbf{\underline{Request Lifecycle:}}

\begin{enumerate}
  \item User submits observation (GPS location) via frontend
  \item Backend validates input (location range lat/lng)
  \item Observation saved ke MongoDB dengan status ``processing''
  \item Python Worker memproses Google Earth Engine — Multi-Source:
    \begin{itemize}
      \item \textbf{Sentinel-1 SAR} (5-15 detik): Query GEE untuk SAR water mask, hitung persentase area basah, deteksi backscatter anomaly
      \item \textbf{Sentinel-2 NDWI} (5-15 detik): Query GEE untuk NDWI, filter cloud cover $<$ 50\%, skip jika terlalu berawan
      \item \textbf{CHIRPS} (3-5 detik): Query data curah hujan 7 hari terakhir, hitung total dan tren
      \item \textbf{OpenLandMap + SRTM} (3-5 detik): Query soil type dan elevation
    \end{itemize}
  \item Semua hasil satelit dikirim ke Gemini 2.0 Flash:
    \begin{itemize}
      \item Input: Data satelit terstruktur (SAR, NDWI, CHIRPS, soil, elevation) + metadata
      \item Gemini menganalisis, memberikan confidence score, natural language assessment, dan rekomendasi
      \item Update status observasi menjadi ``completed''
    \end{itemize}
  \item Frontend polling (tiap 3 detik) via React Query, menampilkan hasil analisis + peta
\end{enumerate}

\textbf{\underline{Confidence Scoring:}}

\begin{verbatim}
Tidak ada weighted formula. Gemini tentukan confidence (0-100)
berdasarkan SELURUH data satelit:

Input Gemini:
  → SAR: {waterPercentage, backscatterMean, confidence}
  → NDWI: {value, cloudCover, available}
  → Rainfall: {total7day, trend}
  → Soil: {type}
  → Elevation: {meters, terrain}

Gemini pertimbangkan:
  → Apakah SAR dan NDWI saling menguatkan?
  → Apakah rainfall konsisten dengan kondisi basah?
  → Apakah tipe tanah mendukung genangan?
  → Apakah elevasi memungkinkan adanya air?
  → Apakah ada anomali antar sumber?

Output Gemini:
  → confidence: 78
  → verdict: "probable"
  → reasoning: "penjelasan natural language"
  → recommendations: [...]
\end{verbatim}

\textbf{\underline{Water Presence Verdict:}}

\begin{tabularx}{\textwidth}{|l|X|}
  \hline
  \textbf{Confidence} & \textbf{Verdict} \\
  \hline
  $\ge$ 80 & DEFINITIVE (confirmed: multiple satellite sources agree on water presence) \\
  60 - 79 & PROBABLE (likely: strong satellite evidence, minor inconsistencies) \\
  30 - 59 & POSSIBLE (some indicators suggest water, but not conclusive) \\
  $<$ 30 & UNLIKELY (no significant water indicators detected) \\
  \hline
\end{tabularx}

\vspace{0.5cm}

\subsectionbox{D. Error Handling \& Graceful Degradation}

Sistem dirancang untuk tetap berfungsi meskipun ada komponen yang gagal:

\begin{tabularx}{\textwidth}{|l|X|}
  \hline
  \textbf{Skenario} & \textbf{Response Sistem} \\
  \hline
  GEE query gagal (timeout/error) & Retry 2x, jika tetap gagal tampilkan ``Satellite data unavailable'' dengan rekomendasi manual survey \\
  Sentinel-1 SAR unavailable & Fallback ke Sentinel-2 NDWI (jika ada) atau tampilkan partial data \\
  Sentinel-2 NDWI unavailable (awan) & Analisis tetap jalan dengan SAR + CHIRPS + soil + elevation \\
  Gemini gagal & Data satelit tetap ditampilkan mentah, tanpa natural language summary \\
  CHIRPS/Soil/Elevation unavailable & Skip source tersebut, analisis jalan dengan data yang ada \\
  GEE query timeout & Retry 2x, fallback ke data parsial \\
  \hline
\end{tabularx}

\vspace{0.5cm}

\subsectionbox{E. Accuracy \& Validation Strategy}

\begin{itemize}
  \item \textbf{Expected Accuracy}: F1-score 0.75-0.90 (lebih tinggi dari pendekatan single-source karena pakai SAR multi-sumber)
  \item \textbf{Validation Method}: Bandingkan output sistem vs ground truth (survei manual) di 10+ lokasi berbeda
  \item \textbf{Test Environments}:
    \begin{itemize}
      \item Large water body (lake, sea, wide river) — SAR sangat akurat
      \item Small water body (pond, narrow river) — batas resolusi 10m
      \item Flooded agricultural field — SAR bagus
      \item Urban area — risiko false positive dari bangunan
      \item Dense forest — SAR limited, vegetasi menutupi
      \item Wetland / swamp — SAR good for detection
      \item Dry area — baseline untuk konfirmasi true negative
    \end{itemize}
  \item \textbf{Accuracy Metrics}: Precision, Recall, F1-score
  \item \textbf{Target}: F1 $\ge$ 0.80 dengan error rate $<$ 5\%
  \item \textbf{Limitations}: Urban false positives, dense vegetation occlusion, temporal gap 1-3 hari
  \item \textbf{Transparency}: Tampilkan data satelit mentah ke user + Gemini reasoning yang transparan
\end{itemize}

\vspace{0.5cm}

\subsectionbox{F. Testing \& Deployment Checklist}

\textbf{\underline{Pre-Launch Checklist:}}
\begin{itemize}
  \item Frontend: all pages tested on mobile (iOS/Android), responsive design verified
  \item Backend: all endpoints tested with proper error codes, rate limits configured
  \item Database: indices created, backup schedule configured (daily MongoDB dump + PostgreSQL dump)
  \item APIs: rate limits set (Gemini 50k/month, GEE quotas managed), caching configured
  \item Security: JWT refresh mechanism working, CORS configured
  \item Load test: simulate 100 simultaneous observations, monitor queue length \& response time
  \item Rollback plan: previous Docker image tagged, database snapshots available
\end{itemize}

\textbf{\underline{Launch Day Checklist:}}
\begin{itemize}
  \item Deploy frontend ke Vercel (auto CI/CD dari GitHub)
  \item Deploy backend ke Railway (Docker container)
  \item Monitor Sentry untuk errors real-time
  \item Test end-to-end dengan sample observation di production
  \item Monitor queue length via Bull dashboard
  \item Backup pre-recorded demo jika WiFi gagal
\end{itemize}

\vspace{0.5cm}

\subsectionbox{G. Limitations \& Honest Disclaimers}

\textbf{Apa yang sistem ini TIDAK bisa lakukan:}
\begin{itemize}
  \item \textbf{Deteksi groundwater} -- Only surface water visible to satelit
  \item \textbf{Prediksi masa depan} -- Only analyzes current moment, no forecasting
  \item \textbf{Menggantikan professional surveys} -- Akurasi ~75-90\%, bukan 99\%
  \item \textbf{Membedakan jenis air} -- Saltwater vs freshwater tidak dibedakan
  \item \textbf{Mengukur kedalaman air} -- Only presence/absence
  \item \textbf{Real-time sempurna} -- Satellite data minimal 1-3 hari
  \item \textbf{Deteksi badan air kecil} -- Resolusi SAR 10m, mungkin lewatkan genangan kecil
\end{itemize}

\textbf{Faktor yang mempengaruhi akurasi:}
\begin{itemize}
  \item Urban areas: Bangunan bisa memantulkan sinyal SAR mirip air (false positive)
  \item Vegetation density: Hutan lebat menutupi ground truth
  \item Terrain complexity: Pegunungan menyebabkan bayangan SAR
  \item SAR saturation: Genangan sangat dangkal mungkin tidak terdeteksi
  \item Temporal gap: Data satelit tidak real-time absolut (1-3 hari delay)
\end{itemize}

\vspace{0.5cm}

\subsectionbox{H. Future Enhancements (Post-MVP)}

\textbf{Phase 2:}
\begin{itemize}
  \item User accounts with profile + observation history + social features
  \item API publik untuk third-party integrations
  \item Mobile app native (React Native / Flutter)
  \item Real-time alert system untuk anomaly detection
\end{itemize}

\textbf{Phase 3:}
\begin{itemize}
  \item IoT sensor integration (optional physical sensors untuk ground truth)
  \item ML model trained on ground truth data (custom confidence model)
  \item Government agency integration (BNPB, PUPR, Kementan data sharing)
  \item Offline mode dengan local satellite caching
\end{itemize}

\textbf{Phase 4:}
\begin{itemize}
  \item Forecasting model (predict water presence 7 days ahead)
  \item Thermal imaging analysis (jika tersedia)
  \item AR visualization untuk visualisasi data
  \item Integration with drone imagery untuk higher resolution
\end{itemize}

\vspace{0.5cm}

\subsectionbox{I. Competitive Analysis}

\begin{center}
\small
\begin{tabularx}{\textwidth}{|l|c|c|c|c|}
  \hline
  \textbf{Feature/Kriteria} & \textbf{Kami} & \textbf{Single Satellite} & \textbf{Photo AI Only} & \textbf{Manual Survey} \\
  \hline
  Multi-source & \greencheck (5 sumber) & \redcross (1 sumber) & \redcross & \greencheck \\
  Tembus awan & \greencheck (SAR) & \redcross (optis) & \greencheck & \greencheck \\
  Peta otomatis & \greencheck (choropleth) & \redcross & \redcross & \redcross \\
  AI analysis & \greencheck (Gemini) & \redcross & \greencheck & \redcross \\
  Biaya/observasi & Rp 0 & 1-5M & 0 & 5-10M \\
  Akurasi & 75-90\% & 70-85\% & 40-60\% & 95\%+ \\
  \hline
\end{tabularx}
\end{center}

\vspace{0.5cm}

\subsectionbox{J. Risk Mitigation}

\begin{tabularx}{\textwidth}{|l|l|l|X|}
  \hline
  \textbf{Risk} & \textbf{Like-lihood} & \textbf{Impact} & \textbf{Mitigation} \\
  \hline
  Gemini quota exceeded & Medium & High & Cache results, batch processing, upgrade tier jika perlu \\
  GEE queries slow & Medium & Medium & Cache 7 hari, set 30-sec timeout, fallback ke cached data \\
  BMKG API down & High & Low & Skip gracefully, continue without weather, fallback ke OpenWeather \\
  Satellite no data & Medium & Medium & Show ``no recent satellite data'', use historical if available \\
  Database corruption & Low & High & Daily backups (MongoDB dump + PostgreSQL dump), test recovery monthly \\
  \hline
\end{tabularx}

\vspace{0.5cm}

\subsectionbox{K. References \& Inspiration}

\begin{itemize}
  \item NDWI methodology: McFeeters, S. K. (1996). The use of Normalized Difference Water Index for discrimination of standing water from soil and vegetation
  \item Citizen science frameworks: Bonney, R., et al. (2009). Citizen Science: A Tool for Integrating Studies of Ecosystems
  \item Self-Determination Theory: Ryan, R. M., \& Deci, E. L. (2000). Intrinsic and Extrinsic Motivations
  \item Google Earth Engine: Gorelick et al. (2017). Google Earth Engine: Planetary-scale geospatial analysis for everyone
  \item Sentinel-2 Band Configuration: ESA Copernicus Sentinel-2 documentation
  \item BMKG API: Badan Meteorologi, Klimatologi, dan Geofisika Indonesia - data.bmkg.go.id
  \item OpenStreetMap: Tile layer for map visualization (openstreetmap.org)
\end{itemize}

\vspace{1.5cm}

% ============================================================================
% FOOTER
% ============================================================================
\begin{center}
  \textit{Proposal ini disusun dengan fokus pada aksesibilitas, transparansi, dan dampak sosial untuk mendukung tujuan pembangunan berkelanjutan di Indonesia. Sistem menggabungkan citizen science, AI analyst, dan multi-source satellite remote sensing dalam satu platform terintegrasi yang accessible melalui smartphone.}
\end{center}

\vspace{1cm}

\begin{flushright}
  \textbf{Berkat ILLAHI} \\
  \textbf{UNIDA Gontor} \\
  Mei 2026
\end{flushright}

\end{document}
